% Role-Playing Test Template
% Canonical simple prompt used throughout this study.
% Compile with: xelatex roleplay_test_template.tex
% The English boxes also compile under pdflatex if the multilingual appendix is removed.
\documentclass[11pt]{article}

\usepackage{fontspec}
\usepackage{xeCJK}
\setmainfont{TeX Gyre Termes}
\setsansfont{TeX Gyre Heros}
\setmonofont{TeX Gyre Cursor}
\IfFontExistsTF{Microsoft YaHei}{\setCJKmainfont{Microsoft YaHei}}{%
  \IfFontExistsTF{Noto Sans CJK SC}{\setCJKmainfont{Noto Sans CJK SC}}{%
    \setCJKmainfont{SimSun}}}
\IfFontExistsTF{Noto Sans Khmer}{%
  \newfontfamily\khmerfont{Noto Sans Khmer}%
}{%
  \newcommand{\khmerfont}{}%
}

\usepackage[margin=1in]{geometry}
\usepackage{booktabs}
\usepackage{array}
\usepackage{tabularx}
\usepackage{xcolor}
\usepackage{enumitem}
\usepackage{fancyvrb}
\usepackage{caption}
\usepackage{hyperref}

\hypersetup{
  colorlinks=true,
  linkcolor=black,
  urlcolor=blue!50!black,
  citecolor=black
}

\DefineVerbatimEnvironment{Prompt}{Verbatim}{%
  frame=single,
  framesep=2mm,
  fontsize=\small,
  baselinestretch=1.05,
  xleftmargin=2mm,
  commandchars=\\\{\}%
}

\newcommand{\slot}[1]{\texttt{\{#1\}}}
\newcommand{\subj}{\textit{subject}}

\title{Role-Playing Test Template}
\author{}
\date{}

\begin{document}
\maketitle

\begin{abstract}
We specify the role-playing test used in this study.
The test is a paired multiple-choice evaluation.
A \emph{base} prompt and a \emph{role-play} prompt share the question, the options, and the output constraint.
They differ only in a single expert-persona sentence placed at the start of the role-play prompt.
The subject name is read from the benchmark label and inserted into that sentence.
Accuracy is the fraction of items whose extracted choice matches the gold label under greedy decoding.
\end{abstract}

\section{Test format}
Each item is a four-way multiple-choice question drawn from a subject of MMLU or MMLU-Redux.
The model receives one user message and must return a single choice.
No few-shot examples, chain-of-thought instruction, or system prompt are added.
The chat template of the evaluated model wraps this user message at inference time.

Two prompts are built for every item.
\begin{description}[leftmargin=1.6em,style=nextline]
  \item[Base.] A task instruction, the question, the four options, and an output constraint.
  \item[Role-play.] The same three blocks, preceded by one persona sentence:
  ``You are an expert in the field of \slot{subject}.''
\end{description}
The comparison therefore isolates the effect of the persona prefix.
Question text, option text, option order, and the required output format are held fixed.

\subsection{Slots}
\begin{itemize}[leftmargin=1.4em]
  \item \slot{subject}: the benchmark subject id, with underscores replaced by spaces and each word capitalized.
  Example: \texttt{college\_biology} becomes \texttt{College Biology}.
  \item \slot{question}: the question stem, unchanged.
  \item \slot{options}: four lines, one choice per line.
  The canonical form labels them \texttt{A.}--\texttt{D.}.
  The MMLU preliminary run labels them \texttt{0.}--\texttt{3.} instead; the rest of the template is unchanged.
\end{itemize}

\section{Canonical templates}
\label{sec:templates}

\subsection{Base prompt}
\begin{Prompt}
Answer the following multiple-choice question. Choose the correct option.

Question: {question}

Options:
A. {choice_0}
B. {choice_1}
C. {choice_2}
D. {choice_3}

Please output only the letter of the correct answer (A, B, C, or D). Do not output any other text.
\end{Prompt}

\subsection{Role-play prompt}
\begin{Prompt}
You are an expert in the field of {subject}. Please answer the following question based on your professional knowledge.

Question: {question}

Options:
A. {choice_0}
B. {choice_1}
C. {choice_2}
D. {choice_3}

Please output only the letter of the correct answer (A, B, C, or D). Do not output any other text.
\end{Prompt}

The role-play template has four blocks, in this order.
\begin{enumerate}[leftmargin=1.4em]
  \item \textbf{Persona.} One sentence that assigns an expert role in \slot{subject} and asks the model to use professional knowledge.
  \item \textbf{Question.} The stem, introduced by \texttt{Question:}.
  \item \textbf{Options.} Four labeled choices.
  \item \textbf{Output constraint.} A request to emit only the letter of the correct answer.
\end{enumerate}
Block~1 is the only text that the base prompt does not contain.
Blocks~2--4 are copied verbatim between the two conditions.

\subsection{Instantiated example}
For the subject \texttt{anatomy}, the role-play prompt is instantiated as follows.
The question is illustrative.
\begin{Prompt}
You are an expert in the field of Anatomy. Please answer the following question based on your professional knowledge.

Question: Which chamber of the heart pumps oxygenated blood into the aorta?

Options:
A. Right atrium
B. Right ventricle
C. Left atrium
D. Left ventricle

Please output only the letter of the correct answer (A, B, C, or D). Do not output any other text.
\end{Prompt}
The matched base prompt is identical after the first sentence is removed and replaced by
``Answer the following multiple-choice question. Choose the correct option.''

\section{Index-labeled variant}
The MMLU preliminary study uses the same wording and the same persona sentence, with numeric labels and a numeric output constraint.
\begin{Prompt}
You are an expert in the field of {subject}. Please answer the following question based on your professional knowledge.

Question: {question}

Options:
0. {choice_0}
1. {choice_1}
2. {choice_2}
3. {choice_3}

Please output only the number of the correct answer (0, 1, 2, or 3). Do not output any other text.
\end{Prompt}
All reported MMLU-Redux results and the hypothesis-verification runs use the letter form in Section~\ref{sec:templates}.

\section{Multilingual persona}
The same four-block layout is used in eight languages: English, Chinese, Japanese, Malay, German, French, Khmer, and Spanish.
Only the surface strings change.
\slot{subject}, \slot{question}, and the option contents are not translated; the subject string remains the English display name defined above.
Table~\ref{tab:persona} lists the persona sentence, which replaces Block~1.
The output constraint asks for a letter (\texttt{A}--\texttt{D}) in the letter setting and for an index (\texttt{0}--\texttt{3}) in the index setting.

\begin{table}[t]
  \centering
  \caption{Persona sentence of the role-play prompt. \slot{subject} is the English display name of the benchmark subject.}
  \label{tab:persona}
  \small
  \begin{tabularx}{\linewidth}{@{}l X@{}}
    \toprule
    Language & Persona sentence \\
    \midrule
    English &
      You are an expert in the field of \slot{subject}.
      Please answer the following question based on your professional knowledge. \\
    Chinese &
      您是\slot{subject}领域的专家。请根据您的专业知识回答以下问题。 \\
    Japanese &
      あなたは\slot{subject}分野の専門家です。専門知識に基づいて次の質問に答えてください。 \\
    Malay &
      Anda seorang pakar dalam bidang \slot{subject}.
      Sila jawab soalan berikut berdasarkan pengetahuan profesional anda. \\
    German &
      Sie sind ein Experte auf dem Gebiet \slot{subject}.
      Bitte beantworten Sie die folgende Frage basierend auf Ihrem Fachwissen. \\
    French &
      Vous êtes un expert dans le domaine de \slot{subject}.
      Veuillez répondre à la question suivante en fonction de vos connaissances professionnelles. \\
    Khmer &
      {\khmerfont អ្នកគឺជាអ្នកជំនាញក្នុងវិស័យ\ \slot{subject}។
      សូមឆ្លើយសំណួរខាងក្រោមដោយផ្អែកលើចំណេះដឹងជំនាញរបស់អ្នក។} \\
    Spanish &
      Usted es un experto en el campo de \slot{subject}.
      Por favor, responda la siguiente pregunta basándose en su conocimiento profesional. \\
    \bottomrule
  \end{tabularx}
\end{table}

\section{Scoring}
Decoding is greedy (\texttt{temperature} $= 0$).
A response is correct when the choice extracted from it equals the gold label.
Extraction searches the generated text, after any thinking span is removed, for a standalone choice token: a letter in $\{\mathrm{A},\mathrm{B},\mathrm{C},\mathrm{D}\}$ or an index in $\{0,1,2,3\}$, depending on the template variant.
If no choice token is found, the item is counted as incorrect.
Accuracy on a subject is the number of correct items divided by the number of items in that subject.
The role-play effect on a subject is the difference
\[
  \Delta = \mathrm{Acc}_{\text{role-play}} - \mathrm{Acc}_{\text{base}}.
\]

\end{document}
